%% kapa: exam=1 round=26 year=2015 period=2 form=A subject=accounting subject_ko=회계학 seq=2 number=42
%% source: Q-Net 감정평가사 기출문제 (cst00309) artlSeq=5138001 file=extracted/1008/5138001/제26회_감정평가사 1차_기출/제26회_감정평가사 1차_기출/2교시_감평_A형.pdf
%% conversion: kordoc 4.24.51 -> tools/gen1.py (auto); review=proofread-sonnet; flags=none
\begin{question}{42}
수익에 관한 설명으로 옳지 않은 것은?
\begin{choices}
  \item 수익은 자본참여자의 출자관련 증가분을 제외한 자본의 증가를 수반하는 것으로서 회계기간의 정상적인 활동에서 발생하는 경제적효익의 총유입을 의미한다.
  \item 판매자가 판매대금의 회수를 확실히 할 목적만으로 해당 재화의 법적 소유권을 계속 가지고 있더라도 소유에 따른 중요한 위험과 보상이 이전되었다면 해당 거래를 판매로 보아 수익을 인식한다.
  \item 수익금액은 일반적으로 판매자와 구매자 또는 자산의 사용자 간의 합의에 따라 결정되며, 판매자에 의해 제공된 매매할인 및 수량리베이트를 고려하여 받았거나 받을 대가의 공정가치로 측정한다.
  \item 이미 수익으로 인식한 금액에 대해서는, 추후에 회수가능성이 불확실해지는 경우에도 이미 인식한 수익금액을 조정하지 아니하고, 회수불가능한 금액이나 더 이상 회수가능성이 높다고 볼 수 없는 금액을 비용으로 인식한다.
  \item 대리관계에서 본인을 대신하여 대리인인 기업이 받는 금액은 경제적효익의 총유입에 해당하므로 대리인인 기업의 수익에 해당한다.
\end{choices}
\end{question}
\answer{5}
